\section{Introduction}
Graph neural network ensembles combine predictors that can compensate for one another's errors. Sharing learned graph weights reduces the number of stored parameters, but also changes which predictors training can acquire. The central question is whether those shared members improve the combined prediction. Correct member decisions can still be lost during averaging, so we evaluate member quality, available correct alternatives and the served prediction together.

GNNM maintains four prediction paths through common graph-layer weights. Established BatchEnsemble factors change the input and output maps \citep{wen2020batchensemble}; related factorized ensembles have been studied for tabular prediction \citep{gorishniy2025tabm} and inside GNN layers \citep{kim2023efficient}. Members retain separate hidden states and graph computations. Our contribution is an empirical study of the predictions learned under these sharing constraints, with weight storage and execution cost reported separately.

The central benchmark covers Roman Empire, Amazon Ratings, Minesweeper, Questions and Tolokers \citep{platonov2023critical}, with eight backbones and ten official splits per dataset. Recorded means favor GNNM over four independent ordinary models in 28 of 40 backbone--dataset comparisons. SAGE and both graph-transformer variants favor GNNM on every dataset, while other backbones include losses. This benchmark compares each method's intended training and validation-selection procedure; matched studies examine sharing within a specified recipe.

The studies reveal different sources of better combined prediction. Matched Roman SAGE improves individual members despite lower disagreement. Internal SAGE on selected WikiCS development instead gains through broader correct-alternative coverage despite weaker members and a calibrated-NLL tradeoff. A conditional neighbor-value extension improves its members relative to a freshly acquired shared baseline and loses fewer correct alternatives during pooling. Receiver conditioning accounts for most of that extension's accuracy gain, and the gain over an equally equipped independent ensemble remains uncertain across three seeds. A packed smaller-model reference adds a concrete storage-budget comparison. These results concern graph- and recipe-specific procedures built from known components; selected development evidence does not establish an unused whole-pipeline effect or a new learning principle.
